\scanpage{80}
\mspage{3}
centimètres de \unsure{longueur|longeur}, fendus longitudinalement au couteau pour que le\\
goût de l'aromate pénètre mieux dans l'eau environnante. Egalement, on\\
peut ajouter dans le pot quelques piments frais ou secs entiers, ils con-\\
tribuent à la conservation du Kimchi et lui donnent un goût un peu plus\\
corsé, sans risquer d'avoir un effet irritant comme la poudre de piment.\\
Quant à l'ail (et toujours pour le Kimchi à l'eau), chaque gousse épluchée\\
est coupée en deux (ou trois, si elle est grosse), et mise à infuser dans\\
le pot, avec les autres aromates. Les aromates ne sont écrasées au mortier\\
que pour le Kimchi ``sans eau''.

La quantité de sel et d'aromates à ajouter pour un pot de Kimchi est\\
réglé au goût, qui doit être engageant, ni trop fort ni trop faible. Il\\
vaut mieux en mettre au début trop peu que trop, quitte à en rajouter le\\
lendemain, ce qui n'offre pas de difficulté. Un Kimchi fortement salé ou\\
épicé pourra être délicieux, à condition qu'il soit servi comme condiment\\
(pour accompagner pommes de terre ou céréales notamment), plutôt que comme\\
plat de légumes.

3.\quad Le matériel nécessaire pour la préparation du Kimchi (mis à part les\\
ingrédients dont il vient d'être question) est des plus simples. Le plus\\
important est le pot en terre ou en grès - il est bon d'en avoir un assor-\\
timent de tailles différentes, suivant les quantités de légumes dont on\\
dispose d'un cas à l'autre. En Corée les pots familiaux de cents ou deux\\
cent litres sont pratique courante et ils y sont peu coûteux - pour notre\\
usage des pots entre trois et vingt litres feront sans doute l'affaire, on\\
en trouve dans le commerce sans trop de mal (quitte à faire une commande au\\
besoin), pour des prix variant suivant la taille, de trente à 200 frs.\\
Il faut aussi disposer de couvercles pour les pots, qui ne doivent jamais\\
rester ouverts quand ils sont en service. D$\underset{\unsure{e}}{\overset{-}{\mathrm e}}$s assiettes font généralement\\
l'affaire, et sont préférables à des planchettes en bois, qui favorisent\\
l'apparition de moisissures. Il est bon quand on prépare le Kimchi de\\
choisir un pot de taille approprié, de façon qu'il soit à peu près rempli\\
et non à demi plein seulement, pour éviter un grand contact d'air.

Pour puiser le Kimchi dans le\unsure{x|s} pot, il convient d'utiliser une\\
spatule ou cuillère en bois pour le Kimchi sans eau, une louche (en bois\\
de préférence) pour le Kimchi à l'eau. Cuillère et louche doivent être\\
propres \uline{et sèches} (les essuyer après lavage !), sous peine d'introduire\\
des fermentations parasites dans le pot et de faire ``tourner'' le Kimchi.\\
Ne pas utiliser la même cuillère ou louche pour plusieurs pots de Kimchi,\\
sans l'avoir lavée et séchée après chaque usage. Eviter de prendre du Kim-\\
chi avec la main dans le pot, sauf si elle vient d'être lavée. Quand on a
\mspage{79}
